%%
%% This is file `sample-acmlarge.tex',
%% generated with the docstrip utility.
%%
%% The original source files were:
%%
%% samples.dtx  (with options: `all,journal,bibtex,acmlarge')
%% 
%% IMPORTANT NOTICE:
%% 
%% For the copyright see the source file.
%% 
%% Any modified versions of this file must be renamed
%% with new filenames distinct from sample-acmlarge.tex.
%% 
%% For distribution of the original source see the terms
%% for copying and modification in the file samples.dtx.
%% 
%% This generated file may be distributed as long as the
%% original source files, as listed above, are part of the
%% same distribution. (The sources need not necessarily be
%% in the same archive or directory.)
%%
%%
%% Commands for TeXCount
%TC:macro \cite [option:text,text]
%TC:macro \citep [option:text,text]
%TC:macro \citet [option:text,text]
%TC:envir table 0 1
%TC:envir table* 0 1
%TC:envir tabular [ignore] word
%TC:envir displaymath 0 word
%TC:envir math 0 word
%TC:envir comment 0 0
%%
%% The first command in your LaTeX source must be the \documentclass
%% command.
%%
%% For submission and review of your manuscript please change the
%% command to \documentclass[manuscript, screen, review]{acmart}.
%%
%% When submitting camera ready or to TAPS, please change the command
%% to \documentclass[sigconf]{acmart} or whichever template is required
%% for your publication.
%%
%%
\documentclass[acmsmall]{acmart}

\usepackage{graphicx}
\usepackage{algorithm}
\usepackage{algpseudocode}
\usepackage{subcaption}
\usepackage{amsmath}
\usepackage{upgreek}
\usepackage{multirow}
\usepackage{booktabs}
\usepackage{url}
\usepackage{comment}
\usepackage{xcolor}
\usepackage[utf8]{inputenc}
\usepackage{array}
\usepackage{makecell}
\newcommand\sd[1]{\textbf{\textcolor{blue}{SD: #1}}	}
%%
%% \BibTeX command to typeset BibTeX logo in the docs
\AtBeginDocument{%
  \providecommand\BibTeX{{%
    Bib\TeX}}}

%% Rights management information.  This information is sent to you
%% when you complete the rights form.  These commands have SAMPLE
%% values in them; it is your responsibility as an author to replace
%% the commands and values with those provided to you when you
%% complete the rights form.
\setcopyright{acmlicensed}
\copyrightyear{0}
\acmYear{0}
\acmDOI{XXXXXXX.XXXXXXX}

%%
%% These commands are for a JOURNAL article.
\acmJournal{TIST}
\acmVolume{0}
\acmNumber{0}
\acmArticle{0}
\acmMonth{0}

%%
%% Submission ID.
%% Use this when submitting an article to a sponsored event. You'll
%% receive a unique submission ID from the organizers
%% of the event, and this ID should be used as the parameter to this command.
%%\acmSubmissionID{123-A56-BU3}

%%
%% For managing citations, it is recommended to use bibliography
%% files in BibTeX format.
%%
%% You can then either use BibTeX with the ACM-Reference-Format style,
%% or BibLaTeX with the acmnumeric or acmauthoryear sytles, that include
%% support for advanced citation of software artefact from the
%% biblatex-software package, also separately available on CTAN.
%%
%% Look at the sample-*-biblatex.tex files for templates showcasing
%% the biblatex styles.
%%

%%
%% The majority of ACM publications use numbered citations and
%% references.  The command \citestyle{authoryear} switches to the
%% "author year" style.
%%
%% If you are preparing content for an event
%% sponsored by ACM SIGGRAPH, you must use the "author year" style of
%% citations and references.
%% Uncommenting
%% the next command will enable that style.
%%\citestyle{acmauthoryear}


%%
%% end of the preamble, start of the body of the document source.
\begin{document}

%%
%% The "title" command has an optional parameter,
%% allowing the author to define a "short title" to be used in page headers.
\title{GCT: A Granger-Causal Transformer for Multivariate Traffic Analysis in Smart Villages}

%%
%% The "author" command and its associated commands are used to define
%% the authors and their affiliations.
%% Of note is the shared affiliation of the first two authors, and the
%% "authornote" and "authornotemark" commands
%% used to denote shared contribution to the research.
\author{Alberto Durán-López}
\email{albduranlopez@ugr.es}
\orcid{0009-0005-8995-869X}
\affiliation{%
  \institution{Research Centre for Information and Communication Technologies
(CITIC-UGR), University of Granada}
  \city{Granada}
  \country{Spain}
}

\author{Daniel Bolaños-Martinez}
\email{danibolanos@ugr.es}
\orcid{0000-0003-0207-2908}
\affiliation{%
  \institution{Research Centre for Information and Communication Technologies
(CITIC-UGR), University of Granada}
  \city{Granada}
  \country{Spain}
}

\author{Suparna De}
\affiliation{%
  \institution{School of Computer Science and Electronic Engineering, University of Surrey}
  \city{Surrey}
  \country{UK}}
\email{s.de@surrey.ac.uk}
\orcid{0000-0001-7439-6077}

\author{Maria Bermudez-Edo}
\email{mbe@ugr.es}
\orcid{0000-0002-2028-4755}
\affiliation{%
  \institution{Research Centre for Information and Communication Technologies
(CITIC-UGR) and IATUR institutes, University of Granada}
  \city{Granada}
  \country{Spain}
}

%%
%% By default, the full list of authors will be used in the page
%% headers. Often, this list is too long, and will overlap
%% other information printed in the page headers. This command allows
%% the author to define a more concise list
%% of authors' names for this purpose.
\renewcommand{\shortauthors}{Durán-López et al.}

%%
%% The abstract is a short summary of the work to be presented in the
%% article.

\begin{comment}
\begin{abstract}
Predicting vehicle traffic optimizes transportation management and urban planning, reduces congestion, and supports infrastructure development. In this paper, we utilize vehicle-detection Internet of Things (IoT) devices to capture real-time and long-term traffic data. Moreover, society’s interconnectedness offers an additional fine-grained data source to add information to IoT data. We use Google Trends data as external variables in this work. Integrating such heterogeneous, complex data streams is challenging for traditional machine learning models that struggle to capture the dynamics of traffic patterns, which are influenced by multiple interdependent factors. Large Language Models (LLMs), which utilize transformer architectures, offer advanced techniques for modeling sequential data. Building on these techniques, this paper introduces Granger-Causal Transformer (GCT), a transformer-based architecture for traffic prediction that integrates an LSTM network with a modified multi-head attention mechanism. This mechanism extends Granger causality to the spatio-temporal domain to analyze all causality relations between features consistently, while capturing long-range dependencies and temporal patterns. Before applying GCT, we generate lagged versions of the Google Trends time series to capture {\color{black} lead and lag effects. Tourists usually make searches about their destination weeks before traveling, so peaks in search interest occur earlier than peaks in weekly traffic volume. Using lags aligns the predictors with weekly traffic volume and allows the model to use past searches to predict future traffic.} {\color{black} We semantically validate the Google Trends terms by comparing each term with a reference string describing the study area, using a language model aligned with the data’s linguistic context.} We then apply a dual filtering process consisting of a test of Granger noncausality to remove time series that lack significant causal relationships and a correlation test to discard insignificant associations. This preprocessing step minimizes noise and redundancy. We evaluate our proposed methodology against classical statistical models, deep learning models, {\color{black}large foundation models}, and transformers {\color{black} across two case studies.} {\color{black}The results demonstrate consistently superior performance and generalizability, with GCT achieving $R^2$ improvements between 47\% and 68\% compared to the best performing baselines across both settings, alongside substantial reductions in MAE and MSE.}
\end{abstract}
\end{comment}

\begin{abstract}
Predicting vehicle traffic optimizes transportation management and urban planning. In this paper, {\color{blue}we combine real-time data from vehicle-detection Internet of Things (IoT) devices with external variables from Google Trends.} Integrating such heterogeneous, complex data streams is challenging for traditional machine learning models that struggle to capture the dynamics of traffic patterns, which are influenced by multiple interdependent factors. {\color{blue}To effectively model these complex, interdependent factors, we introduce the} Granger-Causal Transformer (GCT), a transformer-based architecture for traffic prediction that integrates an LSTM network with a modified multi-head attention mechanism. This mechanism extends Granger causality to the spatio-temporal domain to analyze all causality relations between features consistently, while capturing long-range dependencies and temporal patterns. Before applying GCT, we generate lagged versions of the Google Trends time series to capture lead and lag effects. Tourists usually make searches about their destination weeks before traveling, so peaks in search interest occur earlier than peaks in weekly traffic volume. Using lags aligns the predictors with weekly traffic volume and allows the model to use past searches to predict future traffic. We semantically validate the Google Trends terms by comparing each term with a reference string describing the study area, using a language model aligned with the data’s linguistic context. We then apply a dual filtering process {\color{blue}comprising Granger noncausality and correlation tests is applied to minimize noise and redundancy.}. We evaluate our proposed methodology against classical statistical models, deep learning models, large foundation models, and transformers across two case studies. The results demonstrate consistently superior performance and generalizability, with GCT achieving $R^2$ improvements between 47\% and 68\% compared to the best performing baselines across both settings, alongside substantial reductions in MAE and MSE.
\end{abstract}


\begin{CCSXML}
<ccs2012>
   <concept>
       <concept_id>10010147.10010257.10010321</concept_id>
       <concept_desc>Computing methodologies~Machine learning algorithms</concept_desc>
       <concept_significance>500</concept_significance>
       </concept>
   <concept>
       <concept_id>10010147.10010257.10010293.10010294</concept_id>
       <concept_desc>Computing methodologies~Neural networks</concept_desc>
       <concept_significance>300</concept_significance>
       </concept>
   <concept>
       <concept_id>10003033.10003034</concept_id>
       <concept_desc>Networks~Network architectures</concept_desc>
       <concept_significance>300</concept_significance>
       </concept>
   <concept>
       <concept_id>10010147.10010178</concept_id>
       <concept_desc>Computing methodologies~Artificial intelligence</concept_desc>
       <concept_significance>100</concept_significance>
       </concept>
 </ccs2012>
\end{CCSXML}

\ccsdesc[500]{Computing methodologies~Machine learning algorithms}
\ccsdesc[300]{Computing methodologies~Neural networks}
\ccsdesc[300]{Networks~Network architectures}
\ccsdesc[100]{Computing methodologies~Artificial intelligence}

%%
%% Keywords. The author(s) should pick words that accurately describe
%% the work being presented. Separate the keywords with commas.
\keywords{Granger Causality Test, Transformers, IoT, Traffic Forecasting}

%\received{20 February 2007}
%\received[revised]{12 March 2009}
%\received[accepted]{5 June 2009}

%%
%% This command processes the author and affiliation and title
%% information and builds the first part of the formatted document.
\maketitle

\section{Introduction}

 %{\color{red} https://www.scribbr.com/research-paper/research-paper-introduction/}
 
 
 %{\color{red} PARRAFO 1: CONTEXTO: OK. Información de contexto: Las series temporales se usan mucho en ML y para tráfico}  




Time series analysis plays an important role in the Internet of Things (IoT) domain, contributing to various fields such as healthcare \cite{garcia2025eeg,kumar2020time,ahmed2023multivariate}, industry \cite{raparthy35predictive,de2022deep,aseeri2023effective}, and transportation \cite{zeng2022modeling,xing2022estimating,chauhan2024confined}. Recent advancements in Large Language Models (LLMs) have opened new opportunities to enhance time series methodologies across these domains. LLM-inspired architectures emphasize dependency modeling and multivariate sequence processing, transforming tasks such as behavior prediction, anomaly detection, and real-time decision-making. Specifically, in transportation, LLM-based approaches can advance traffic forecasting by leveraging historical data to build more robust predictive models, improving traffic management strategies and infrastructure planning \cite{li2024traffic}.

Advances in traffic prediction have used a range of machine learning {\color{black}(ML)} techniques for multivariate time series data. Initially, statistical models like AutoRegressive Integrated Moving Average (ARIMA) were preferred because they capture temporal dependencies in a linear manner \cite{shahriari2020ensemble, kumar2022time}. However, as traffic patterns became more complex and the volume of data increased, deep learning models such as Recurrent Neural Networks (RNN), Long Short-Term Memory (LSTM), and Gated Recurrent Units (GRU) became widely used \cite{ismaeel2023traffic, liu2020traffic, hussain2021intelligent}. These models can capture complex nonlinear relationships and time-based dependencies that ARIMA cannot efficiently handle. Recently, transformer-based architectures have gained popularity due to their ability to manage long-term dependencies through attention mechanisms \cite{reza2022multi, yu2023multi,bolanos2025sasd}. Advancements in large language models (LLMs), which are built upon transformer architectures, have demonstrated their effectiveness in handling sequential data, which can 
% Although primarily utilized in natural language processing, the techniques and mechanisms developed for LLMs can inspire and
enhance methodologies in time series analysis, including in applications such as traffic forecasting \cite{reza2022multi, yu2023multi}.


With the move towards sustainable tourism, especially in rural regions, demand forecasting and optimized resource allocation require real-time as well as historical and fine-grained data about the local conditions. Though IoT devices are a primary source of data, complexities associated with installation ease (due to local geography), device maintenance, data storage and communication, mean that secondary data sources can serve to augment spatio-temporal data about rural regions. 
% is important in order to achieve , which involves  optimizing  and strategic planning. 
Social data, particularly Google Trends, have emerged as a promising additional data source that can mitigate reporting delays and data incompleteness, providing real-time insights to predict future trends \cite{choi2012predicting}. Incorporating Google Trends data into mixed-data sampling models has been shown to significantly enhance forecast accuracy \cite{bangwayo2015can} and it is widely used to predict tourism traffic demand \cite{xiao2020tourism, onder2017forecasting}.  However, Google Trends data may be lagged relative to other sources; for example, a tourist might search for a destination a few weeks or even months before traveling. This anticipatory search behavior introduces an additional challenge in aligning the timing of these search metrics with actual travel events.


%\sd{we need a short para here to introduce here the importance/usefulness of having social data, i.e. Google trends data in our case, as an additional data source + the challenges it brings.}

A key aspect of neural network (NN) modeling is understanding that not all input variables contribute equally to the forecasting task. To address this, studies have introduced techniques into ML models to evaluate and prioritize the importance of variables. Tank et al. \cite{tank2021neural} proposed component-wise MLP (cMLP) and component-wise LSTM (cLSTM), which incorporate causality into deep learning models to automatically select and prioritize input series with significant causal influence on the output series. Others, such as \cite{wang2024rethinking}, introduce Causal Attention Mechanism (CATT), a framework that uses causal inference to remove spurious correlations by separating valid (causal) features from invalid (non-causal) features. Similarly, CausalFormer \cite{kong2024causalformer} introduced a transformer-based model architecture that integrates causality with attention by applying causal convolutions before the attention mechanism and using penalties to identify and prioritize causal input series, without relying on statistical tests for causal validation.


%CausalFormer enforces a strict temporal priority constraint, assuming that the order of causation is fixed, and it does not explicitly model sequential temporal dependencies. 

%{\color{red} PÁRRAFO 3: GAP en el Estado del rate Problema con el estado del arte: No se usa la causalidad de manera X. Aquí pon el problema que tienes ahora mismo en el estado del arte.}
%hecho nuevo en un nuevo párrafo sacando el gap de la literatura.

Approaches such as CATT and CausalFormer introduce causal structures by modifying attention mechanisms to emphasize or isolate features that are causal based on observed patterns. While these strategies help mitigate spurious correlations, the resulting causal inferences are driven primarily by architectural heuristics rather than by statistical verification. Consequently, there remains an opportunity for methods that systematically integrate and validate causal relationships within the attention framework through formal empirical tests.


%{\color{red} PÁRRAFO 4: Describir nuestra solución Nuestra solución: To address this GAP... metemos una matriz de causalidad en el attention, y además un LSTM y además..}

This gap motivates the development of our Granger-Causal Transformer (GCT), an architecture designed to explicitly incorporate Granger causality into the attention mechanism to statistically test and model causal dependencies between pairs of variables or inputs. Granger causality represents the causality between two time series as a regression task, with a ``Granger'' cause determined if one time series can successfully predict another. {\color{black}Here, Granger causality is used as a forecasting criterion: past values of a variable $X$ improve predictions of another variable $Y$ when we condition on the lagged values of $Y$. This does not imply structural or true causation \cite{granger1969investigating}.}
GCT combines an LSTM layer with a modified multi-head attention mechanism. The LSTM layer captures sequential dependencies, while the attention layer integrates a causality mask, which accounts for Granger-causal relationships between variables, ensuring that the model captures both temporal dependencies and causal interactions, while also reflecting the strength and precedence of relationships between variables. Prior to GCT prediction, we validate the external variables using {\color{black}a language model (LM) aligned with the data’s linguistic context}, and then apply a Granger causality test as a feature selection filter, combined with a correlation filter. This dual-filtering approach ensures that only causally relevant and highly correlated time series are selected and subsequently passed to the GCT architecture for prediction. Additionally, we conduct an ablation test to evaluate the contribution and significance of each element of our methodology. 


%We validate our proposed methodology using a case study of smart villages in the Alpujarra region of the Sierra Nevada mountains in southern Spain, which serves to evaluate the proposed method for analyzing vehicle traffic data in rural areas. {\color{black}In this study, we forecast weekly vehicle counts for the whole study area, the three-village area connected by a single access road (Pampaneira, Bubión, and Capileira). We target a medium-term forecasting horizon and focus on whole-area predictions rather than road-segment or per-village forecasts, reflecting local mobility patterns and aiming to generalize to region-level demand beyond this specific case.}


We validate the methodology in {\color{black} two case studies}. In the first study (Alpujarra, southern Spain), we forecast weekly vehicle counts for the three-village area (Pampaneira, Bubión, and Capileira) connected by a single access road. {\color{black} In the second study (Salt Lake City, USA), we aggregate weekly counts from four stations positioned along the main roads leading into the city (from north, south, east, and west). In both cases, we forecast several weeks in advance, as tourists typically plan their trips well in advance. The predictions are made at the destination-wide level rather than for specific roads or locations, as tourists usually visit multiple places within the area.}


%Main contributions
The main contributions of this work are highlighted as follows:
\begin{itemize}   
    \item We present GCT, a transformer-based architecture that modifies the attention mechanism to adjust attention weights based on causality between series. In particular, we incorporate a mask that introduces the Granger causality test into the attention mechanism.
    \item We introduce a feature validation step {\color{black}using a language model} as a validator, ensuring that only semantically relevant features are retained for a subsequent dual statistical filtering process, which applies both a Granger causality test and a correlation filter to refine the selection of lagged features for time series prediction.
    \item We conduct an ablation study to evaluate our method by exploring how the inclusion or exclusion of the causality mask in the transformer, and the two feature selection mechanisms-causality and correlation filtering-affect model performance. This analysis helps assess the contribution of each component to the overall prediction metrics.
    \item {\color{black}We validate the robustness and generalizability of our GCT pipeline across two distinct contexts: a rural setting with Spanish-language data and an urban setting with English-language data, demonstrating its effectiveness in different environments.}
\end{itemize}


%{\color{red} PÁRRAFO 6: Estructura del artículo}

The remainder of the paper is organized as follows. Section \ref{sec:related_work} reviews the related work. Section \ref{sec:preliminaries} discusses the theoretical foundations of Granger causality and the transformer attention mechanism. Section \ref{sec:methods} presents our methodology. Section \ref{sec:case_study} outlines the {\color{black}initial case study}. Section \ref{subsec:experiments} presents the experimental results, including the ablation study (Section \ref{subsec:ablation_test}) {\color{black}and a second case study to validate generalizability (Section \ref{subsec:urban_repro}). This section also discusses the findings (Section \ref{sec:discussion}) and details limitations and future work (Section \ref{sec:limitations})}. Finally, Section \ref{sec:conclusions} concludes the paper.


%Additionally, attention mechanisms are important for capturing long-range dependencies and nuanced temporal relationships in multivariate traffic data \cite{alharthi2024enhanced,altunay2023hybrid,lin2024bn,oyando2023rnn,tang2023swinlstm}. Although a few architectures integrate attention \cite{dong2023time,lv2024transportation}, they still do not implement causality analysis, which limits their ability to fully leverage the richness of multivariate data.

%Statistical models such as Autoregressive Integrated Moving Average (ARIMA) are employed in traffic forecasting \cite{shahriari2020ensemble,liu2020traffic}. However, ARIMA faces difficulties when dealing with nonlinearity or long-term dependencies \cite{albeladi2023time,pirani2022comparative}, and these limitations become even more pronounced in real-world traffic datasets, which are not only noisy but also often multivariate, containing multiple interdependent variables that statistical models cannot easily handle. To address these challenges, some researchers use deep learning techniques, which are adept at handling nonlinear relationships and long-term dependencies \cite{khan2023short,li2023novel}. 
%Other studies use time series models, such as Gated Recurrent Units (GRU) or Long Short-Term Memory (LSTM), dealing with multivariate datasets. \cite{ray2023arima,habib2024wifog,hu2024dynamic}. When dealing with time series data, effective feature selection must account for temporal interdependencies between variables \cite{verdonck2024special}, utilizing methods such as Pearson's and Spearman's correlation coefficients to evaluate linear and monotonic relationships, respectively \cite{shiekh2022comparison,kathirgamanathan2021correlation}. Another approach is to use causality analysis for feature selection, which has been explored to better capture temporal interdependencies between variables \cite{pan2023high,yu2020causality}.


%Although causality and correlation have been explored separately in the literature, their combination in these architectures has not been studied. Hence, we propose an algorithm that uses both causality and correlation for feature selection in multivariate time series. 

%La revista pone por puntos la contribución: https://www.sciencedirect.com/journal/advanced-engineering-informatics

\section{Related Work} \label{sec:related_work}


%Artículos que predicen el tráfico utilizando modelos estadísticos o ML. Es un párrafo general. Hay al menos un artículo para cada algoritmo que se usa en la comparativa de baselines.
Classical statistical models like ARIMA \cite{kumar2022time, shahriari2020ensemble} and its multivariate variant ARIMAX \cite{su2020daily} are widely utilized for predicting traffic flow and analyzing time series data in public transportation systems. ARIMAX has proven more effective than ARIMA by incorporating external factors, as shown in studies focusing on travel times \cite{novak2023data} and accident predictions \cite{ihueze2018road}. {\color{black}This concept of using exogenous data has been extended to sources like social media, where Twitter activity and sentiment from the previous night have been used to predict the next morning's commute traffic \cite{yao2021twitter}.} With advancements in deep learning, models that capture complex traffic patterns have emerged. RNNs, particularly LSTM and GRU, have enhanced traffic flow prediction by modeling temporal dependencies \cite{ismaeel2023traffic, narmadha2023spatio, sroczynski2023road}. {\color{black}Architectural refinements, such as using stacked bidirectional LSTMs, have been proposed to capture both forward and backward temporal dependencies more effectively, even addressing missing data issues \cite{cui2020stacked}.} Additionally, combining convolutional layers (e.g., Conv1D) with LSTMs has improved the ability to capture spatial-temporal patterns \cite{long2021air}, {\color{black}a concept also extended to architectures utilizing GRU-based layers \cite{huang2025prediction}}. Recently, transformer-based models have gained importance for their capability to handle large-scale datasets and capture long-range dependencies, as evidenced by their application in spatio-temporal traffic prediction \cite{tian2023multi} and traffic congestion forecasting \cite{ramana2023vision}.





%{\color{black}These lines of work highlight two complementary paths: sparsity penalties that enforce structure during training and attention models that recover structure by post hoc interpretation.}


% 1. Se empieza explicando de forma GENERAL causalidad.
%2. Se introducen diferentes formas de introducir causalidad en los modelos:
    
    %En este primer párrafo:
    %long2021air usa Granger
    %cMLP/cLSTM usan LASSO (group, sparse group or hierarchical)

    %En el párrafo de abajo (2 abajo)
    %wang2024rethinking aplica un "causal inference framework" that removes influence of confounders.
    %Causalformer usa transformers+causality. Transformer architecture used for temporal causal discovery;    
{\color{black}Several} studies have incorporated causality to determine the importance of variables in multivariate models. For instance, Long et al. \cite{long2021air} improve traffic prediction models by utilizing causality to identify the most influential variables within neural networks. Other researchers have combined causality with ML models to enhance prediction performance. {\color{black}One approach to this combination is through regularization penalties.} For example, Tank et al. \cite{tank2021neural} introduce the component-wise Multilayer Perceptron (cMLP) and component-wise Long Short-Term Memory (cLSTM), neural network architectures designed for Granger causality detection in multivariate temporal series. cMLP employs separate multilayer perceptrons for each target series, applying group lasso, sparse group lasso, or hierarchical lasso penalties to the input weights. This penalization helps identify causal relationships between the input series and the output series and also enables automatic lag selection. Similarly, cLSTM utilizes individual LSTM networks for each output series, with group lasso penalties imposed on the input weights to determine causal influences. Causality between time series is then determined by the presence of non-zero weights connecting the input series to the output within these networks. 
%{\color{black}The key advantage of these lasso-based penalties is that they induce sparse connections, automatically detecting which input series and time lags are causally relevant to each output. This sparsity not only improves interpretability but also stabilizes estimation when dealing with many series and lags.} %These models rely on structured sparsity and regularization to identify causal relationships. 

%comento el siguiente párrafo ya que esto NO va en el related work.
%{\color{black}In contrast to this embedded approach (which uses penalties), our work adopts a filter-based method: we apply explicit Granger causality and correlation tests prior to the model to select inputs, which are then encoded into a causal attention mask. This filter-based approach separates variable selection from model training, improving interpretability and avoiding the need for penalty hyperparameter tuning.}. 

%Esto usa Granger como feature selection, lo pongo más abajo mejor.
%{\color{black}More recent frameworks have also employed Granger causality as a feature selection method for multitask LSTM models, enhancing the predictions by also introducing theory-informed loss functions \cite{fafoutellis2025theory}.}

%Seguimos con la idea del párrafo anterior, pero ahora usando transformers, es decir, incorporar causality+transformers.
{\color{black}Rather than relying on sparsity-inducing penalties, other researchers have explored incorporating causality directly into model architectures \cite{kong2024causalformer} or through explicit causal inference mechanisms \cite{wang2024rethinking}, such as, encoding causal relationships within transformer architectures themselves}.
For instance, Wang and Zhou \cite{wang2024rethinking} reinterpret the role of attention mechanisms in transformers from a causality perspective, concentrating specifically on image data and not adapting their approach to domains such as time-series forecasting. They introduce the concept of confounders, which are variables that influence both the input data and the attention mechanism, leading to spurious correlations that can negatively impact model performance. {\color{black}To address confounding without regularization penalties,} they propose a causal inference framework that removes the influence of confounders by applying interventions based on the do-operation and front-door criterion, enabling the attention mechanism to focus on truly causal relationships rather than spurious correlations. {\color{black}Taking a different architectural approach,} Kong et al. \cite{kong2024causalformer} introduce Causalformer, a transformer architecture used for temporal causal discovery in time series data. Causalformer employs multi-kernel causal convolution blocks to learn causal representations by convolving each time series with causal kernels under temporal priority constraints. After training the transformer, they apply Regression Relevance Propagation (RRP) as a post hoc interpretability method to decompose the model’s predictions into relevance scores. This post-attention analysis enables the extraction and construction of meaningful causal relationships from the trained model. 

%Se añaden Large Foundation Models. Se explican Lag Llama, Chronos, TimeGPT, TimesFM y Moirai, pero sólo se incluyen los dos últimos como baselines.
%Estos modelos SÍ usan la arquitectura Transformer, pero NO implementan la causalidad explícitamente en el mecanismo de atención, que es la contribución clave de GCT. 
{\color{black}Recent popular examples that utilize transformers include large foundation models for time series, pre-trained on web-scale data and used in zero-shot settings. These models vary in their design and accessibility. Some, such as Lag-Llama \cite{rasul2023lag} and Chronos \cite{ansari2024chronos}, are designed for probabilistic forecasting, outputting predictive distributions rather than single point forecasts. Other models, such as TimeGPT \cite{garza2023timegpt}, are primarily accessible via proprietary APIs. In contrast, publicly available models, such as TimesFM \cite{das2024decoder} and Moirai \cite{woo2024unified}, focus on deterministic point-forecasting. Among these, TimesFM is a decoder-only transformer, pre-trained primarily on Google Trends and Wikipedia page views, which aligns well with demand and seasonality patterns. It produces strong zero-shot forecasts. Because it does not take exogenous time series at forecast time, external effects are captured only indirectly through recurring patterns learned during pretraining and visible in the series history. Another large foundational model is Moirai, which is a unified transformer trained on diverse time series corpora; it supports multivariate inputs and exogenous covariates at inference, enables multi-horizon forecasting, and can operate in zero-shot mode using official checkpoints. Both provide competitive baselines, though neither explicitly encodes statistical causality nor structured lag selection within attention.}



% Filtros ANTES de la arquitectura.
Regarding feature selection prior to the architecture, various approaches are employed. For example, Li et al. \cite{li2022pearson} use the Pearson Correlation Coefficient in a Broad Learning System (BLS) to enhance prediction accuracy by filtering highly correlated features, while Salman et al. \cite{salman2022stability} employ Spearman's Rank Correlation Coefficient to improve feature selection stability in high-dimensional datasets, using methods like the arithmetic mean and L2 norm. In contrast, some authors have leveraged causality for feature selection \cite{yu2020causality}. Long et al. \cite{long2021air} utilize causal relationships to select features prior to their hybrid model combining 1D convolutional neural networks and LSTM networks. {\color{black}In a similar way, Fafoutellis and Vlahogianni \cite{fafoutellis2025theory} apply a Granger causality test to select a reduced, causally-informed input space for a multitask LSTM.} Yu et al. \cite{ling2024fair} introduce the Fair Causal Feature Selection Algorithm (FairCFS), an approach that accounts for causal relationships to ensure fairness in classification tasks.  %{\color{black}Compared with penalty-based embedded selection, our filter-based selection combines causality and correlation tests outside the forecasting model. This keeps the forecasting objective simple, improves reproducibility, and separates variable screening from parameter shrinkage.}

Considering the absence of methods that directly incorporate Granger causality into the attention mechanism of transformer models, and the lack of approaches that combine both correlation and causality-based feature selection within such architectures, we propose our methodology to include these components. 



\section{Background} \label{sec:preliminaries}


\subsection{Granger Causality Test}
The Granger Causality test \cite{granger1969investigating}, is a statistical method used to determine if one time series can predict another. The main idea is that if a time series $X$ granger-causes another time series $Y$, it means that changes in $X$ happen first and are reflected later in $Y$, implying there is a delay between the two. In this case, past values of $X$ provide information that helps predict $Y$ beyond the information in the past values of $Y$ alone \cite{papaspyropoulos2024validity}.

Consider two time series, $X_t$ and $Y_t$. The autoregressive models used in the Granger causality test are the Restricted Model (without $X$) and the Full Model (with $X$). The Restricted Model, which uses past values of $Y$ to predict $Y$, is defined as follows:
\begin{equation}
    Y_t = \alpha_0 + \alpha_1 Y_{t-1} + \alpha_2 Y_{t-2} + \cdots + \alpha_p Y_{t-p} + \epsilon_t.
\end{equation}

The Full Model, which uses past values of $Y$ and $X$ to predict $Y$, is defined as follows:


\begin{equation}
    Y_t = \beta_0 + \beta_1 Y_{t-1} + \beta_2 Y_{t-2} + \cdots + \beta_p Y_{t-p} + \gamma_1 X_{t-1}  + \gamma_2 X_{t-2} + \cdots + \gamma_p X_{t-p} + \nu_t.
\end{equation}

The null hypothesis of the Granger causality test is that $X_t$ does not Granger-cause $Y_t$, meaning that all the coefficients $\gamma_j$ are equal to zero. This hypothesis is tested using an F-test, which compares the sum of squared residuals (SSR) from the restricted and unrestricted models:

\begin{equation}
    F = \frac{(\text{SSR}_{\text{restricted}} - \text{SSR}_{\text{full}}) / p}{\text{SSR}_{\text{full}} / (n - 2p - 1)},
\end{equation}

where:
\begin{itemize}
    \item $\text{SSR}_{\text{restricted}}$ is the sum of squared residuals of the model without $X$.
    \item $\text{SSR}_{\text{full}}$ is the sum of squared residuals of the model with $X$.
    \item $p$ is the number of lags (number of previous periods used to predict the current value).
    \item $n$ is the number of observations.
\end{itemize}

The F-test statistic calculated above follows a specific probability distribution known as the F-distribution, which depends on the number of lags and the number of observations. To make a decision, we compare the calculated F-test value to a critical value obtained from the F-distribution table at a chosen significance level of 5\%, a commonly used threshold in statistical hypothesis testing to assess result significance \cite{bruns2019lag}. If the F-test value is greater than the critical value, we reject the null hypothesis and conclude that $X_t$ Granger-causes $Y_t$. This means that $X_t$ contains useful information for predicting $Y_t$ that is not already captured by the past values of $Y_t$ alone.



\subsection{Transformer: Attention Mechanism}

%Self-attention 

Transformers are neural network architectures designed to leverage parallel processing and effectively capture long-range dependencies within data \cite{vaswani2017attention}. Central to their design is the attention mechanism, which allows the model to dynamically focus on the most relevant parts of the input sequence when generating each element of the output, enabling it to model contextual relationships across the sequence.  A fundamental component of the attention mechanism is self-attention, where the input sequence itself serves as the queries, keys, and values. Given an input sequence $X = [x_1, x_2, \dots, x_n]$, where $x_i$ is the feature vector for the $i$-th element, the self-attention mechanism is computed as follows:

\begin{equation}
    \text{Attention}(Q, K, V) = \text{softmax}\left(\frac{QK^T}{\sqrt{d_k}}\right)V
\end{equation}

In this formulation, the input $X$ is linearly transformed into three distinct matrices: queries ($Q$), keys ($K$), and values ($V$). These are computed as:

\begin{equation}
    Q = XW_Q, \quad K = XW_K, \quad V = XW_V
\end{equation}
where $W_Q$, $W_K$, and $W_V$ are learnable weight matrices. The dot product of the queries and keys, $QK^T$, results in attention scores that are scaled by $\frac{1}{\sqrt{d_k}}$, where $d_k$ is the dimension of the keys, to avoid excessively large values. In a multivariate NN applied to time series data, the matrix $\frac{QK^T}{\sqrt{d_k}}$ represents {\color{black}pairwise similarities between different time steps in these input series, rather than only similarities across series}. Each element in this matrix corresponds to the similarity between the query from one {\color{black} time step} and the key from another {\color{black} time step}, allowing the model to capture relationships between the inputs and adjust the attention scores accordingly \cite{lin2017structured}. The softmax function is then applied to normalize these scores, ensuring they sum to 1. The resulting weights are then used to compute a weighted sum of the values ($V$), producing the final output of the self-attention mechanism.
In the multi-head attention mechanism, the process of self-attention is repeated multiple times in parallel. Instead of computing a single set of attention scores, the input is projected into several different subspaces, producing multiple sets of queries, keys, and values ($Q_i$, $K_i$, $V_i$) for each head $i$.

%The resulting matrix now reflects both the original attention scores and the influence of Granger causality, allowing the model to better capture causal dependencies in the data.

\section{Methodology} \label{sec:methods}

Our methodology (see Fig. \ref{fig:methodology}) begins with the validation of input series using {\color{black}a language model aligned with the data’s linguistic context}, followed by the generation of lagged versions of the time series, which are filtered to retain only those series that exhibit significant causal relationships and correlations, as detailed in Subsection \ref{subsec:featureSelection} and step 1 of Algorithm \ref{alg:pseudo}. The filtered features are then normalized and passed to our architecture to generate predictions. Our architecture consists of an LSTM layer to capture temporal dependencies {\color{black}(see Subsection \ref{subsec:LSTM} and step 2 of Algorithm \ref{alg:pseudo})} and a multi-head attention layer with a causality mask to integrate causal relationships (see Subsection \ref{subsec:transformer} and step 3 of Algorithm \ref{alg:pseudo}). Additionally, the model includes a normalization layer to stabilize learning, residual connections to facilitate gradient flow and prevent vanishing gradients \cite{yu2019review}, and dense layers with dropout to learn abstract representations and reduce overfitting \cite{srivastava2014dropout}. This methodology provides a structured temporal and causality-aware approach for multivariate time-series forecasting. {\color{black}Finally, Subsection \ref{subsec:modeltrain} details the training process, hyperparameters, and the cross-validation strategy used for evaluation.}
%The red boxes in Figure \ref{fig:methodology} highlight the components used in the ablation test, as detailed in subsection \ref{subsec:ablation_test}.


\begin{figure}[h]
\centering
	\includegraphics[width = 0.9\linewidth]{Figures/BlockDiagram.pdf}
	\caption{Methodology}
        \Description{Image of the followed methodology}
        \label{fig:methodology}
\end{figure}



\subsection{Feature Selection}
\label{subsec:featureSelection}

{\color{black}We use external social data from search engines commonly used by visitors to plan their trips. Specifically, for our use cases, we use Google Trends, which provides data on how often a given term was searched on Google each week \footnote{\url{https://trends.google.com/trends/}}. In Google Search, a higher volume of a term within a week raises the relative search-interest index reported by Google Trends. We treat these indices as leading indicators. People typically search for destinations and activities before visiting. Hence, sustained increases in search interest are expected to precede subsequent tourist inflows. Importantly, these searches typically occur pre-trip from the tourists' regions of origin, where internet access is generally widespread, mitigating concerns about low internet penetration within the destination itself. 
}

{\color{black}
 To obtain external series for a destination (study area), we begin with an initial set of social data terms, for example, in our use case, Google Trends terms, that represent the local context and then expand them with terms frequently searched for the same region and topic.
}

Then, we establish semantic relationships to ensure that only semantically relevant time series are considered. {\color{black} Because aligning the language model with the linguistic and contextual domain of the data is important for semantic validation, our methodology adapts this step based on the Google Trends terms' language. For example, in our first case study (Section \ref{sec:case_study}), which uses Spanish queries, we use BETO, a pre-trained Spanish language model\footnote{https://huggingface.co/dccuchile/bert-base-spanish-wwm-uncased} \cite{canete2023spanish}.} BETO follows the BERT-Base architecture and applies the Whole Word Masking (WWM) technique for improved token representations \cite{devlin2018bert}. The model was trained on a large Spanish corpus, using a vocabulary of approximately 31,000 subword units generated with the SentencePiece algorithm \cite{kudo2018sentencepiece}, and optimized over 2 million training steps. BETO generates dense vector embeddings that effectively capture both semantic and syntactic relationships between terms. {\color{black}In contrast, for our English-language generalizability study (Section \ref{subsec:urban_repro}), we use the BERT model\footnote{https://huggingface.co/bert-base-uncased} to ensure proper alignment.}

{\color{black}Using the selected language model, we encode a reference string that describes the study area and each candidate term (query terms in the search engine)}. We then compute cosine similarity (see Eq. \ref{eq:cosine_sim}) between the reference string embedding and each term embedding and retain only terms with similarity $\ge \tau$. By applying a cosine similarity threshold to discriminate relevant semantic relationships \cite{reimers2019sentence}, we filter out semantically unrelated terms and retain only the most pertinent ones. {\color{black}Next, for those Google Trends terms, we obtain their associated weekly time series that represents the number of times each term was searched in Google. This yields a semantically aligned set of terms whose weekly numeric values serve as the external time series.}

\begin{equation} \label{eq:cosine_sim}
  \text{cosine\_similarity}(A, B) = \frac{A \cdot B}{|A| |B|}
\end{equation}

After that, we account for temporal dependencies and incorporate historical information by generating lagged versions of each series. This allows the model to consider past values when making predictions, enhancing its ability to capture dynamic relationships within the data. Furthermore, incorporating irrelevant time series can lead to overfitting, increased computational complexity, and diminished forecasting performance \cite{hyndman2018forecasting}. To address these challenges, we apply causal relationship filtering and correlation-based filtering during preprocessing. 


%DECIR POR QUÉ TENGO QUE HACER EL LAG. Normalm
%como la relación entre esta X y la Y está desplazada en el tiempo (porque personas buscan con antelación, pero no sabes cuánto desplazamiento), hago muchos y luego aplico los filtros. explicarlo desde un punto de vista general.


The next step involves creating lagged versions of the social data. This step ensures that the model accounts for historical values that may influence the target variable $Y$ over different time periods. From the generated series, we select those that have a significant relevance.
%causal relationship with the target variable, as determined by the Granger causality test.
Let $\{X_t\}$ denote a time series. The lagged versions of this series are represented as $X_{t-1}, X_{t-2}, \dots, X_{t-p}$, where $p$ is the maximum lag considered. Specifically, we select the series where the p-value of the Granger test satisfies $p \leq 0.05$ \cite{bruns2019lag}, indicating that they contribute valuable predictive information. The selection process described in step 1 of Algorithm \ref{alg:pseudo} can be summarised as follows:
\begin{equation} \label{eq:equation_causality}
    \text{Select } X_{t-i} \text{ such that } \text{p-value of Granger test}(X_{t-i} \to Y) \leq 0.05
\end{equation}


Since including all the series and their respective lags may introduce redundancy and noise, we select only the most relevant ones based on their correlation. Only series with a correlation coefficient greater than a predefined threshold, $\alpha$, are retained:
\begin{equation} \label{eq:equation_correlation}
|\text{Corr}(X_{t-i}, Y_t)| > \alpha 
\end{equation}
{\color{black}Here, $\alpha \in [0,1]$ is a threshold applied to the absolute Pearson correlation between each lagged candidate time series and the target. Lower $\alpha$ makes the filter less restrictive and retains more predictors, increasing the risk of redundancy and overfitting, whereas higher $\alpha$ retains fewer, more strongly correlated time series, at the risk of discarding informative signals.} 

This filtering guarantees that the retained time series not only have a causal impact on $Y$ but also exhibit a strong linear relationship with it. By enforcing both causality and correlation, we ensure that the model incorporates only those predictors that provide relevant predictive information while reducing complexity by eliminating less relevant series. {\color{black}In our methodology, these Granger and correlation screens serve only as a filtering step for the lagged social time series. The retained signals are then modelled by the non-linear architecture (see Section \ref{subsec:transformer}), which captures both linear and non-linear interactions.} {\color{black} Unlike penalty-based approaches such as Lasso, our filter-based method separates variable screening from model training. This avoids penalty hyperparameter tuning, improves transparency, and allows the transformer to focus on capturing complex interactions among the validated causal predictors through the attention mask.}


{\color{black}
\subsection{LSTM Component} \label{subsec:LSTM}

The first layer of our architecture is an LSTM layer, which functions as a sequential encoder. LSTMs are designed to process information sequentially, which makes them well-suited for capturing local temporal dependencies and sequential order \cite{sak2014long}. Its role is to first process the input features, injecting a temporal inductive bias and transforming the sequence of inputs into a sequence of enriched, context-aware hidden state vectors. This denoises the inputs and highlights relevant temporal patterns. This feature extraction provides the subsequent GCT layer with a cleaner, contextually richer input.

This component involves two sequential operations. First, the input features are passed through an LSTM layer, which is configured with recurrent units to produce a sequence of vectors known as hidden states that encode and summarize the learned temporal dependencies. Second, this sequence of hidden states is processed by a Layer Normalization operation. The purpose of this normalization step is to stabilize the distribution of the hidden state activations. This is an important step to mitigate potential vanishing or exploding gradient issues and ensure that the inputs fed to the subsequent GCT layer are well-scaled, which facilitates a more stable and efficient training process \cite{ba2016layer}.

}

\subsection{{\color{blue}GCT architecture}} \label{subsec:transformer}

GCT incorporates an attention mechanism that accounts for Granger-causal relationships between time series in a transformer. To that end, we construct a causality matrix $C$ which serves as a mask. This mask is applied to the attention mechanism to adjust the attention similarity weight matrix. Fig. \ref{fig:attention_mechanisms} shows both the regular transformer and the proposed attention mechanism with causality mask, which is formulated as follows:
\begin{equation} \label{eq:equation_attention}
    \text{Attention}(Q,K,V) = \text{softmax}\left(\frac{QK^T}{\sqrt{d_k}} \odot (1 + C)\right)V
\end{equation}
We construct the matrix $C$ by evaluating pairwise Granger causality between time series $i$ and $j$, where:
\begin{equation}
C_{ij} = \frac{1}{\text{best\_lag}(i,j)} \cdot \mathbb{I}(X_i \rightarrow X_j)
\end{equation}
Here, the indicator function $\mathbb{I}(X_i \rightarrow X_j)$ equals 1 if series $X_i$ Granger-causes series $X_j$, and 0 otherwise. If series $i$ does not Granger-cause series $j$, then $C_{ij} = 0$, meaning no adjustment is made to the attention score for this pair. However, if series $i$ significantly Granger-causes series $j$, the corresponding entry $C_{ij}$ is set to $\frac{1}{\text{best\_lag}(i,j)}$, where $\text{best\_lag}(i,j)$ is {\color{black} an integer in $\{1,\dots,p_{\max}\}$, representing} the shortest significant time delay found between the two series. {\color{black}$p_{\max}$ represents the maximum viable lag, beyond which the relationship between the time series becomes too weak or noisy to be informative. For example, in our use cases, we set $p_{\max}=12$, based on previous literature that states that many pre-trip searches occur up to 12 weeks before travel \cite{claude2020predicting,bangwayo2015can}. Larger $p_{\max}$ values (e.g., $16$--$20$) are hard to estimate on weekly data and tend to dilute the information and add noise. Smaller $p_{\max}$ values (e.g., $6$--$10$) can miss useful longer lags.} The indicator function ensures that only causal relationships contribute to the matrix $C$. This is reinforced by the observation that shorter lags typically indicate stronger causal relationships, as they suggest that the influence of one series on another is immediate \cite{gunther2022reconstruction}. The use of the inverse lag in the entry $C_{ij}$ reflects this notion, where a shorter lag corresponds to a higher value of $C_{ij}$, emphasizing the immediacy and strength of the causal effect \cite{peng2020dli}. The term $1 + C$ is then multiplied element-wise (Hadamard product) with the similarity matrix $\frac{QK^T}{\sqrt{d_k}}$, meaning that the attention score $\frac{q_i k_j^T}{\sqrt{d_k}}$ between query $q_i$ and key $k_j$ is adjusted individually based on the causal relationship $C_{ij}$. This localized adjustment allows the attention mechanism to focus on specific causal interactions without affecting the entire similarity matrix.

After the element-wise multiplication, the softmax function normalizes the adjusted similarity matrix into a valid probability distribution. This process, detailed in step 3 of Algorithm \ref{alg:pseudo}, preserves the original attention mechanism's ability to capture contextual relationships while incorporating causal information. As a result, the model learns both contextual and temporal causal relationships, providing a richer, causality-aware representation of the data.

\begin{figure}[H]
    \centering
\includegraphics[width=0.45\linewidth]{Figures/new layer.pdf}
    \caption{Comparison of attention mechanisms. The left image shows the original attention mechanism, while the right image shows the attention mechanism with causality mask.}
    \label{fig:attention_mechanisms}
\end{figure}

\begin{algorithm}[H]
\caption{GCT Training and Inference}
\begin{algorithmic}[1]
\Require $Y$: Target Variable, $X = \{X_1, X_2, \dots, X_n\}$: Predictor Variables, $p$: Number of lags
\Ensure $\hat{Y}$: Predicted Target Variable

\State $X_{{\color{black}LM}} \leftarrow \{ x_i \in X \mid \text{cosine\_similarity}(\text{{\color{black}LM}}(x_i), \text{{\color{black}LM}}({\color{black}reference\_string})) \geq \text{threshold} \}$
\State $Z \leftarrow [Y; X_{{\color{black}LM}}]$
\State $Z_{\text{norm}} \leftarrow \text{Normalize}(Z)$




\vspace{0.2cm}
\State \textbf{Step 1}: {\color{black}Feature Selection:} Causality and Correlation
\For{$i \leftarrow 1$ \textbf{to} $n$}
    \For{$j \leftarrow 1$ \textbf{to} $p$}
        \State Generate lagged series: $X_{i,t-j}$
        \State Retain $X_{i,t-j}$ if it satisfies causality and correlation thresholds
    \EndFor
\EndFor

\vspace{0.2cm}
\State \textbf{Step 2}: LSTM {\color{black}Component}
\State $H_{\text{LSTM}} \leftarrow \text{LSTM}(Z_{\text{norm}})$ \Comment{Extract temporal dependencies}
\State $H_{\text{norm}} \leftarrow \text{LayerNorm}(H_{\text{LSTM}})$

\vspace{0.2cm}
\State \textbf{Step 3}: {\color{blue}GCT Architecture}
\State Initialize causality matrix: $C \leftarrow \mathbf{0}^{(n+1) \times (n+1)}$

\For{$i \leftarrow 1$ \textbf{to} $n+1$} \Comment{$Y$ is included as $X_{n+1}$}
    \For{$j \leftarrow 1$ \textbf{to} $n+1$}
        \State $C_{ij} \leftarrow \frac{1}{\text{best\_lag}(X_i, X_j)} \cdot \mathbb{I}{(X_i \rightarrow X_j)}$
    \EndFor
\EndFor

\State $Q \leftarrow H_{\text{norm}} W_Q$, $K \leftarrow H_{\text{norm}} W_K$, $V \leftarrow H_{\text{norm}} W_V$
\State $\text{Attention}(Q,K,V) \leftarrow \text{softmax}(\frac{Q K^T}{\sqrt{d_k}} \odot (1 + C)) V$
\State $H_{\text{attention}} \leftarrow \text{Attention}(Q,K,V)$

\vspace{0.2cm}
%\State \textbf{Step 4}: Position-wise Feed Forward and Prediction
\State Add residual connection: $H_{\text{residual}} \leftarrow H_{\text{norm}} + H_{\text{attention}}$
\State Concatenate outputs: $H_{\text{concat}} \leftarrow \text{Concat}(H_{\text{residual}}, H_{\text{norm}})$
\State $H_{\text{FFN}} \leftarrow \text{FeedForwardBlock}(H_{\text{concat}})$ \Comment{Dense + Dropout layers}
\State $\hat{Y} \leftarrow \text{Dense}(H_{\text{FFN}}, \text{units}=1)$

\Return $\hat{Y}$

\end{algorithmic}
\label{alg:pseudo}
\end{algorithm}














\subsection{Model training and evaluation} \label{subsec:modeltrain}

For model training, we employed the Adam optimizer with a learning rate of 1e-4 and a batch size of 64. We used Python 3.10.12 for programming. We present the complete layer structure of our model in Table \ref{tab:neural_model}. Alongside each layer, we provide the best values of hyperparameters. 


\begin{table}[H]
\centering
\resizebox{\textwidth}{!}{
\begin{tabular}{@{}cccc@{}}
\toprule
\textbf{Layer No} & \textbf{Type}                                       & \textbf{Shape}                                                    & \textbf{Parameters}                         \\ \midrule
0                 & Input                                         & (N\textsubscript{Batch}, num\_variables)                                  & N\textsubscript{Batch} = 64                \\
1                 & LSTM                                                & (N\textsubscript{Batch}, lstm\_units)                              & lstm\_units=32                                   \\
2                 & LayerNormalization                                  & (N\textsubscript{Batch}, lstm\_units)                              & -                                                \\
3                 & Multi-HeadAttention with Causality Mask            & (N\textsubscript{Batch}, num\_variables)                              & num\_heads=4, key\_dim=num\_variables                \\
4                 & Add                                                 & (N\textsubscript{Batch}, num\_variables)                              & -                                                \\
5                 & Concatenate                                   & (N\textsubscript{Batch}, lstm\_units + num\_variables)                & -                                                     \\
6                 & Dense                                               & (N\textsubscript{Batch}, dense\_units)                             & dense\_units=64                                  \\
7                 & Dropout                                             & (N\textsubscript{Batch}, dense\_units)                             & dropout=0.15                                     \\
8                 & Dense                                               & (N\textsubscript{Batch}, dense\_units//2)                          & -                                                \\
9                 & Dropout                                             & (N\textsubscript{Batch}, dense\_units//2)                          & -                                                \\
10                & Dense                                               & (N\textsubscript{Batch}, 1)                                        & -                            \\ \bottomrule
\end{tabular}
}
\caption{GCT architecture}
\label{tab:neural_model}
\end{table}


To ensure a fair and robust comparison, we applied 10-fold Cross-Validation (CV) using the TimeSeriesSplit method, which evaluates model performance across different temporal segments of the data \cite{berrar2019cross}. The cross-validation process starts with 50\% of the data for training and 20\% for testing, incrementing the training size gradually up to 80\%, while maintaining the test size fixed at 20\%. Fig. \ref{fig:cv} shows the distribution of the training and testing sets across these 10 CV. To evaluate the model performance, we use the metrics: Mean Absolute Error (MAE), Mean Squared Error (MSE), and the coefficient of determination ($R^2$). Among these, $R^2$ serves as the primary basis for comparison, as it best captures the explanatory power of the models in predicting the target variable \cite{chicco2021coefficient}.

\begin{figure}[ht]
\centering
    \includegraphics[width = 0.8\linewidth]{Figures/cv.pdf}
    \caption{TimeSeriesSplit schema}
    \label{fig:cv}
\end{figure}





\newpage 
\section{Case Study}\label{sec:case_study}


This section presents our case study conducted in the Alpujarra region, located in southern Spain, which consists of three villages: Pampaneira, Capileira, and Bubión. This case study serves to evaluate the proposed method for predicting the weekly vehicle counts in the area. 


The data collection process includes setting up the Automatic Number Plate Recognition (ANPR) devices, positioned as shown in Fig. \ref{fig:map}. We used four Hikvision IP cameras equipped with ANPR systems based on Deep Learning, see Fig. \ref{fig:detection}. These devices have vehicle detection sensors, a 2MP resolution, 2.8-12mm varifocal lenses, and IR LEDs with a range of 50 meters. The ANPR cameras were positioned at the entry and exit points of each village to cover and monitor all vehicular traffic entering the area \cite{bolanos2024spatio, duran2025route}.






\begin{figure}[h]
\centering
	\includegraphics[width = 0.8\linewidth]{Figures/map_V2.pdf}
	\caption{Distribution of the sensors (in red points) in the study area}
        \label{fig:map}
\end{figure}
 
\begin{figure}[h]
\centering
	\includegraphics[width = 0.8\linewidth]{Figures/detection.pdf}
	\caption{Automatic License Plate Recognition with IoT camera devices}
        \label{fig:detection}
\end{figure}
 

%\subsection{Data Description} \label{subsec:dataset}

The ANPR cameras recorded the license plates and timestamps of each vehicle passing through their coverage area. The dataset spans 24 months (from February 2022 to January 2024). After collecting the raw data, which includes the number plate, camera identifier, date, and direction (see Table \ref{tab:raw_dataset}), {\color{black}we applied a data cleaning to recover missing digits in the plate number detections using the Damerau-Levenshtein distance algorithm \cite{duran2024routerecoverer}. Next, we aggregated the cleaned vehicle detections into weekly vehicle counts for the study area, which serve as the target variable $Y$ in our forecasting model.} 



\begin{table}[ht]
\centering
\footnotesize
\resizebox{0.6\linewidth}{!}{%
\begin{tabular}{lll}
\toprule
\textbf{Variable} & \textbf{Description} & \textbf{Source} \\ 
\midrule
num\_plate & License plate number & ANPR Cameras \\ 
camera\_ID & Identifier for ANPR camera & ANPR Cameras \\ 
date & Timestamp of appearance & ANPR Cameras \\ 
direction & Direction of vehicle (entry or exit) & ANPR Cameras \\ 
\bottomrule
\end{tabular}%
}
\caption{Raw Dataset}
\label{tab:raw_dataset}
\end{table}

%In addition to the vehicle count data, we enriched the dataset with relevant search trends obtained from the \textit{Google Trends} platform\footnote{\url{https://trends.google.es/trends/}}. Specifically, we queried traffic-related keywords and phrases relevant to the study area, such as searches for nearby villages and rural tourism terms available on the platform. We identified these keywords as potential indicators of traffic activity in the region. We collected the search data using the PyTrends API \footnote{\url{https://pypi.org/project/pytrends/}}, which retrieves relative Google Trends data and organized it into a weekly time series format. These weekly search trends serve as predictor variables $X = \{X_1, X_2, \dots, X_{15}\}$ in our model, representing a total of 15 distinct weekly search queries, as shown in Table \ref{tab:trends_dataset}.

%dejo el párrafo anterior para que veas los cambios... era bastantes cambios y he tenido que hacerlo así y no modificando palabras del anterior.

In addition to the vehicle count data, we enriched the dataset with relevant {\color{black} Google Trends time series} obtained from the Google Trends platform\footnote{\url{https://trends.google.es/trends/}}. {\color{black}To that end, we began with an initial set of Google Trends terms describing the case study (see Table \ref{tab:trends_dataset}). We then expanded this set and obtained the associated weekly time series using the PyTrends API\footnote{\url{https://pypi.org/project/pytrends/}}.} These weekly time series serve as predictor variables $X = \{X_1, X_2, \dots, X_{15}\}$ in our model, representing a total of 15 distinct weekly search queries, as shown in Table \ref{tab:trends_dataset}.




\begin{table}[ht]
\centering
\footnotesize
\resizebox{\linewidth}{!}{
\begin{tabular}{p{2cm}p{2,5cm}m{7cm}p{2,3cm}}
\toprule
\textbf{Variable} & \multicolumn{2}{c}{\textbf{Description}} & \textbf{Source} \\
\midrule
date           & \multicolumn{2}{l}{Weekly dates from Feb 2022 to Jan 2024.} & ANPR Cameras \\
vehicle\_count & \multicolumn{2}{l}{Weekly vehicle counts.}                   & ANPR Cameras \\
\cmidrule(lr){2-3} & \textbf{Group} & \textbf{{\color{black}Google Trends terms}} & \\
\midrule
\multirow{2}{*}{google\_trends} 
                  & {\color{black}Initial set of Google Trends terms} 
                  & \parbox[t]{7cm}{Pampaneira, Capileira, Bubión\\
                     (\textit{These refer to the main villages of the  {\color{black}case study}).}} 
                  & {\color{black}Manually selected} \\ 
  
                & {\color{black}Expanded set of Google Trends terms} 
                & \parbox[t]{7cm}{Trevélez, Trevélez Spain, Lanjarón Spain, Soportújar Spain, Órgiva Spain, Alpujarra, the Alpujarra granadina, Granada Alpujarra, villages of the Alpujarra, Granada Alpujarra villages, Alpujarra Magna, rural house Granada.\\
                  (\textit{These terms denote the broader {\color{black}geographical context of the study area, capturing related locations and regional searches focused on the entire area}).}} 
                & PyTrends API \\
\bottomrule
\end{tabular}
}
\caption{{\color{black}Description of the Rural} Case Study dataset}
\label{tab:trends_dataset}
\end{table}




{\color{black} As an example, Figure \ref{fig:examples} plots the Google Trends time series of the initial set of Google Trends terms from Table \ref{tab:trends_dataset}, alongside the weekly vehicle counts. The blue arrows highlight typical lead and lag patterns, where peaks in search interest tend to precede increases in vehicles by a few weeks, with seasonal variation. This initial set and the expanded set of terms serve as inputs to our methodology, where they are semantically validated with the BETO language model using the reference string ``Turismo rural en la Alpujarra'' (see Equation \ref{eq:cosine_sim}). We  then generate lagged candidates and retain only those that pass both the Granger causality test and the correlation test (see Equations 
\ref{eq:equation_causality} and \ref{eq:equation_correlation}). The resulting validated, lagged time series serve as the predictor variables for the GCT {\color{blue}architecture}.
}




\begin{figure}[H]
\centering
	\includegraphics[width = 0.8\linewidth]{Figures/examples.pdf}
	\caption{{\color{black} Weekly Google Trends time series for Pampaneira, Capileira, and Bubión (left y-axis), together with vehicles per week entering the area (right y-axis). Dashed annotations indicate temporal shifts (lag) between peaks in relative search interest in the Pampaneira and the subsequent increase in vehicle counts in the area.}}

    \label{fig:examples}
\end{figure}







\section{Results and Discussion}  \label{subsec:experiments}

{\color{black}
This section evaluates the proposed GCT methodology. We first present descriptive statistics of the dataset to establish the baseline characteristics of the traffic patterns in the study region. Table \ref{tab:traffic_basic} reports traffic indicators at daily, weekly, and monthly scales for the Alpujarra region, where the total resident population was 1,209 according to the Spanish National Statistics Institute (INE)\footnote{\url{https://www.ine.es}}. Although data are collected at multiple temporal scales, we use weekly aggregations to align with Google Trends time series. %All figures refer to the observation window from January 2022 to January 2024, covering 104 weeks and 1,000,658 vehicle counts.}

\begin{table}[ht]
\centering
\footnotesize
\setlength{\tabcolsep}{4pt}
\resizebox{\textwidth}{!}{
\begin{tabular}{l*{12}{r}}
\toprule
& \multicolumn{4}{c}{Daily} & \multicolumn{4}{c}{Weekly} & \multicolumn{4}{c}{Monthly} \\
\cmidrule(lr){2-5} \cmidrule(lr){6-9} \cmidrule(lr){10-13}
& All region & Pampaneira & Bubión & Capileira & All region & Pampaneira & Bubión & Capileira & All region & Pampaneira & Bubión & Capileira \\
\midrule
Mean       & 1349.3 & 1322.0 & 1256.1 & 1191.8 &  9444.8 &  9253.9 &  8792.8 &  8342.6 & 39290.6 & 38496.3 & 36578.1 & 34705.0 \\
Median     & 1270 & 1244 & 1179 & 1115 &  9446 &  9258 &  8820 &  8386 & 40473 & 39586 & 37616 & 35901 \\
P25        &  975 &  952 &  912 &  867 &  8112 &  7969 &  7506 &  7137 & 36102 & 35382 & 33646 & 31939 \\
P75        & 1659 & 1647 & 1554 & 1457 & 10464 & 10259 &  9848 &  9491 & 45704 & 45057 & 43194 & 40875 \\
Min        &  333 &  329 &  307 &  283 &  3459 &  3400 &  3177 &  3054 &  36857 & 34100 & 31177 & 30254 \\
Max        & 3357 & 3269 & 3076 & 2887 & 15985 & 15474 & 14814 & 13968 & 61135 & 59645 & 56723 & 53742 \\
Std.\ dev. &  492.1 &  483.3 &  459.5 &  433.8 &  2206.3 &  2157.9 &  2055.8 &  1947.3 & 11107.8 & 10861.1 & 10292.6 &  9707.1 \\
\bottomrule
\end{tabular}
}
\caption{{\color{black}Traffic volume metrics by temporal scale and location. The table reports descriptive statistics for the full study area (Alpujarra region) and for each town (Pampaneira, Bubión, Capileira) at daily, weekly, and monthly aggregation levels.}}
\label{tab:traffic_basic}
\end{table}

Following this descriptive analysis, we conduct a series of experiments to validate the GCT {\color{blue}architecture}. We begin in Subsection \ref{sec:sensitivity_anal} with a sensitivity analysis to justify our selection of preprocessing hyperparameters. Subsection \ref{sec:results} presents the main experimental results for our rural case study, comparing GCT against all baselines. We then conduct an ablation study in Subsection \ref{subsec:ablation_test} to analyze the individual contributions of our model's key components. To validate the model's robustness, we evaluate its performance in a new setting in Subsection \ref{subsec:urban_repro}. Finally, we synthesize all findings in Subsection \ref{sec:discussion}.
}







{\color{black}
\subsection{Sensitivity analysis} \label{sec:sensitivity_anal}

We conduct a sensitivity analysis for the case study to assess how the semantic filter threshold $\tau$ (Equation \ref{eq:cosine_sim}) and the correlation threshold $\alpha$ (Equation \ref{eq:equation_correlation}) affect both predictive performance and computational cost. Starting from the 15 original Google Trends terms in Table \ref{tab:trends_dataset}, we vary $\tau\in\{0.25, 0.50, 0.75\}$ and keep only those terms whose cosine similarity with the study-area reference string is at least $\tau$. For every semantically validated term, we obtain its weekly Google Trends time series, generate up to $p_{\max}=12$ weekly lags, and test each lagged time series for Granger causality against the target series $Y_t$ (Equation \ref{eq:equation_causality}). The column Pass Granger $p$-value $<0.05$ in Table \ref{tab:sensitivity_tau_alpha_final} reports the number of lagged series that are significant. We then apply the correlation screen $|\mathrm{Corr}(X_{t-i},Y_t)|>\alpha$ with $\alpha\in\{0.25, 0.50, 0.75\}$ to those same Google Trends time series; the column Final time series counts the Google Trends time series that pass both filters.
}


\begin{table}[H]
\centering
\footnotesize
\resizebox{\textwidth}{!}{%
\begin{tabular}{cccccccccc}
\toprule
\makecell{\textbf{$\tau$}\\\textbf{(Eq. \ref{eq:cosine_sim})}} &
\makecell{\textbf{$\alpha$}\\\textbf{(Eq. \ref{eq:equation_correlation})}} &
\makecell{\textbf{Semantically}\\\textbf{validated terms}\\\textbf{(Eq. \ref{eq:cosine_sim})}} &
\makecell{\textbf{Lagged}\\\textbf{series}\\\textbf{($p_{\max}=12$)}} &
\makecell{\textbf{Pass Granger}\\\textbf{$p$-value $<0.05$}\\\textbf{(Eq. \ref{eq:equation_causality})}} &
\makecell{\textbf{Final time series}\\\textbf{Correlation filter}\\\textbf{(Eq. \ref{eq:equation_correlation})}} &
\textbf{MAE ($\pm$SD)} & \textbf{MSE ($\pm$SD)} & \textbf{$R^2$ ($\pm$SD)} & \textbf{Time (s) ($\pm$SD)} \\
\midrule
0.25 & 0.25 & 15 & 180 & 48 & 26 & 0.0490 $\pm$ 0.0150 & 0.0040 $\pm$ 0.0021 & 0.704 $\pm$ 0.192 & 29.35 $\pm$ 0.96 \\
0.25 & 0.50 & 15 & 180 & 48 & 20 & 0.0460 $\pm$ 0.0146 & 0.0035 $\pm$ 0.0019 & 0.742 $\pm$ 0.165 & 27.85 $\pm$ 0.80 \\
0.25 & 0.75 & 15 & 180 & 48 & 13 & 0.0558 $\pm$ 0.0163 & 0.0053 $\pm$ 0.0023 & 0.615 $\pm$ 0.240 & 25.90 $\pm$ 0.70 \\
\midrule
0.50 & 0.25 & 12 & 144 & 40 & 22 & \textbf{0.0433 $\pm$ 0.0159} & 0.0037 $\pm$ 0.0024 & 0.760 $\pm$ 0.155 & 28.15 $\pm$ 0.84 \\
\textbf{0.50} & \textbf{0.50} & 12 & 144 & 40 & 20 & 0.0438 $\pm$ 0.0143 & \textbf{0.0032 $\pm$ 0.0019} & \textbf{0.773 $\pm$ 0.132} & 27.37 $\pm$ 0.71 \\
0.50 & 0.75 & 12 & 144 & 40 & 14 & 0.0451 $\pm$ 0.0148 & 0.0034 $\pm$ 0.0020 & 0.751 $\pm$ 0.140 & 26.45 $\pm$ 0.71 \\
\midrule
0.75 & 0.25 & 8 & 96 & 24 & 14 & 0.0516 $\pm$ 0.0155 & 0.0045 $\pm$ 0.0022 & 0.665 $\pm$ 0.195 & 27.05 $\pm$ 0.77 \\
0.75 & 0.50 & 8 & 96 & 24 & 12 & 0.0469 $\pm$ 0.0149 & 0.0037 $\pm$ 0.0020 & 0.735 $\pm$ 0.150 & 26.20 $\pm$ 0.69 \\
0.75 & 0.75 & 8 & 96 & 24 & 8 & 0.0604 $\pm$ 0.0167 & 0.0061 $\pm$ 0.0025 & 0.546 $\pm$ 0.260 & \textbf{24.80} $\pm$ 0.62 \\
\bottomrule
\end{tabular}%
}
\caption{{\color{black}Sensitivity analysis over the cosine-similarity threshold $\tau$ and the correlation threshold $\alpha$, starting from 15 original Google Trends terms (Table \ref{tab:trends_dataset}). Semantically validated terms are Google Trends terms retained by the semantic validation step via cosine similarity with the reference string using the selected language model (Eq. \ref{eq:cosine_sim}). Lagged series are the number of semantically validated terms multiplied by $p_{\max}$, with $p_{\max}=12$. Pass Granger ($p$-value $<0.05$) counts lagged series with significant Granger causality (Eq.~\ref{eq:equation_causality}). Final time series are those that pass the final correlation filter (Eq.~\ref{eq:equation_correlation}). Metrics are mean $\pm$ SD over the 10 folds.}}
\label{tab:sensitivity_tau_alpha_final}
\end{table}
{\color{black} As shown in Table \ref{tab:sensitivity_tau_alpha_final}, the configuration of $\tau=0.5$ and $\alpha=0.5$ yielded the best predictive performance (highest $R^2$ and lowest MSE). Furthermore, the semantic similarity threshold ($\tau=0.5$) is consistent with its use in related literature \cite{leung2022semantic} and the correlation threshold ($\alpha=0.5$) is also consistent with conventional cutoffs for strong linear association \cite{adler2010quantifying}. Hence, we fixed $\tau=0.5$ and $\alpha=0.5$ for our experiments.
}

\subsection{Results} \label{sec:results}
For evaluation, we compared our methodology with various baseline algorithms in both univariate and multivariate settings using cross-validation. In the univariate case, we implemented ARIMA, GRU, LSTM, LSTM with a multi-head attention layer {\color{black}and TimesFM \cite{das2024decoder}}. For the multivariate setting, we included ARIMAX (the multivariate extension of ARIMA), Multivariate GRU, Multivariate LSTM, Multivariate Conv1D with LSTM \cite{long2021air}, cMLP and cLSTM \cite{tank2021neural}, Causalformer \cite{kong2024causalformer} {\color{black}and Moirai \cite{woo2024unified}}. In the univariate cases, the method uses only the historical values of the target variable $Y$. In the multivariate cases, the method uses past values of the target variable $Y$ along with Google Trends time series data as external predictor variables.

The results are presented in Table \ref{tab:sota} and Fig. \ref{fig:sota}. The statistical models (ARIMA, ARIMAX) showed lower $R^2$ scores, of $-0.1552$ and $0.1029$, respectively. Univariate models, such as GRU, LSTM, and LSTM with Attention, demonstrated moderate performance, with average $R^2$ scores ranging from $0.3847$ to $0.4596$. The multivariate GRU and LSTM models showed slight improvements over their univariate counterparts but did not reach the performance levels of our proposed model. {\color{black}Large foundation models showed mixed results. Moirai achieved an $R^2$ of $0.0896$, while TimesFM performed better ($R^2=0.2542$), which is expected given its pre-training on Google Trends data. However, both foundation models fell significantly short compared to causality-aware approaches, such as CausalFormer, and GCT. However, they yield faster execution times (1.70s for TimesFM, 5.01s for Moirai), as they lack explicit causal relationship 
modeling.} The cLSTM model performed poorly, even worse than the statistical models. cMLP performed better than cLSTM but still did not achieve good $R^2$ scores. Causalformer showed slightly higher $R^2$ values compared to previous methods but did not exceed the 0.5 $R^2$ threshold. Our model achieved the lowest MAE of $0.0438$, the lowest MSE of $0.0032$, and the highest $R^2$ score of $0.7734$, indicating a better fit to the data compared with the other evaluated models. These results represent improvements of $30\%$ in MAE, $56\%$ in MSE, and $68\%$ in $R^2$ compared to the univariate LSTM with attention, the second-best model. Execution times were longer for our model,  averaging $27.37$, compared to the other algorithms, which required at most $23.63$ seconds. However, this additional computation time is compensated by the consistently superior performance of our model across all metrics, as observed in nearly all CV folds, as shown in Table \ref{tab:sota} and Fig. \ref{fig:sota}.


\begin{table}[ht]
\centering
\footnotesize
\resizebox{\textwidth}{!}{%
\begin{tabular}{l|c|c|c|c}
\toprule
 & \multicolumn{4}{c}{\textbf{Metric ($\pm$ Std Dev)}} \\ 
\cmidrule(lr){2-5}
\textbf{Model} & \textbf{MAE} & \textbf{MSE} & \textbf{R\textsuperscript{2}} & \textbf{Time (s)} \\ 
\midrule

ARIMA 
& 0.0982 $\pm$ 0.0143 
& 0.0155 $\pm$ 0.0036 
& -0.1552 $\pm$ 0.2342 
& 0.8981 $\pm$ 0.1022 \\ 

ARIMAX 
& 0.0865 $\pm$ 0.0125 
& 0.0119 $\pm$ 0.0030 
& 0.1029 $\pm$ 0.2484 
& \textbf{0.4096 $\pm$ 0.0415} \\ 

Univariate GRU
& 0.0712 $\pm$ 0.0258 
& 0.0085 $\pm$ 0.0049 
& 0.3847 $\pm$ 0.3085 
& 19.0396 $\pm$ 0.7736 \\

Multivariate GRU
& 0.0910 $\pm$ 0.0164 
& 0.0141 $\pm$ 0.0068 
& -0.1300 $\pm$ 0.7257 
& 21.9710 $\pm$ 0.2746 \\

Univariate LSTM
& 0.0684 $\pm$ 0.0134 
& 0.0084 $\pm$ 0.0036 
& 0.3922 $\pm$ 0.2332 
& 18.9211 $\pm$ 0.2085 \\

Univariate LSTM with Attention
& 0.0627 $\pm$ 0.0217 
& 0.0073 $\pm$ 0.0051 
& 0.4596 $\pm$ 0.3423 
& 23.6318 $\pm$ 1.0318 \\

Multivariate LSTM
& 0.0866 $\pm$ 0.0223 
& 0.0122 $\pm$ 0.0096 
& 0.0280 $\pm$ 0.7968 
& 21.4400 $\pm$ 0.6348 \\

Multivariate Conv1d+LSTM \cite{long2021air}
& 0.0718 $\pm$ 0.0168 
& 0.0083 $\pm$ 0.0035 
& 0.3555 $\pm$ 0.3284 
& 18.9420 $\pm$ 0.3810 \\


{\color{black}TimesFM \cite{das2024decoder}}
& 0.0810 $\pm$ 0.0086
& 0.0107 $\pm$ 0.0026
& 0.2542 $\pm$ 0.1188
& 1.70 $\pm$ 0.1969 \\

{\color{black}Moirai \cite{woo2024unified}}
& 0.1067 $\pm$ 0.0246
& 0.0183 $\pm$ 0.0070
& 0.0896 $\pm$ 0.1963
& 5.01 $\pm$ 0.2232 \\

cLSTM \cite{tank2021neural}
& 0.1400 $\pm$ 0.0347 
& 0.0464 $\pm$ 0.0205 
& -0.7696 $\pm$ 0.9802 
& 10.33 $\pm$ 0.6386 \\

cMLP \cite{tank2021neural}
& 0.0955 $\pm$ 0.0157 
& 0.0161 $\pm$ 0.0048 
& 0.0455 $\pm$ 0.1703 
& 1.61 $\pm$ 0.6030 \\

Causalformer \cite{kong2024causalformer}
& 0.0783 $\pm$ 0.0084 
& 0.0101 $\pm$ 0.0018 
& 0.2891 $\pm$ 0.1300 
& 8.49 $\pm$ 0.4633 \\

GCT (Our model)
& \textbf{0.0438 $\pm$ 0.0143}
& \textbf{0.0032 $\pm$ 0.0019}
& \textbf{0.7734 $\pm$ 0.1323}
& 27.37 $\pm$ 0.7053 \\

\bottomrule
\end{tabular}%
}
\caption{Cross-Validation Performance - Methods with Execution Times {\color{black}- Rural Case Study}}
\label{tab:sota}
\end{table}



\begin{figure}[H]
\centering
    \includegraphics[width = \linewidth]{Figures/r2_boxplot_vertical.pdf}
    \caption{Comparison of $R^2$ results}
    \label{fig:sota}
\end{figure}




\subsection{Ablation test} \label{subsec:ablation_test}

We analyzed the proposed method with different configurations to explore how the inclusion or exclusion of causality filtering (Equation \ref{eq:equation_causality}), correlation filtering (Equation \ref{eq:equation_correlation}), and the causal attention mask (Equation \ref{eq:equation_attention}) affects the model's performance. These components correspond to the red boxes in Fig. \ref{fig:methodology}. This analysis allowed us to quantify the effect of each component both individually and in combination. The full model corresponds to our proposed methodology, incorporating causality filtering, correlation filtering, and the causality mask in the GCT architecture. To evaluate the impact of these components, we progressively removed each feature, ultimately reaching a base model that lacks all three components.



%\ref{tab:results_ablation} shows the results for all CV methods, and \ref{fig:results_ablation} displays a boxplot illustrating these results. \ref{fig:train_loss_ablation,fig:validation_loss_ablation} presents the training loss and validation loss of these methods across all epochs during the last CV fold. %Additionally, \ref{fig:radar_ablation} shows the normalized metrics to compare each model’s performance. The radar plot (\ref{fig:radar_ablation}) is structured so that the best-performing model appears fully filled, representing the best metric values, while the worst-performing model remains close to the center, indicating minimal performance. 

We can see in Table \ref{tab:results_ablation} and Fig. \ref{fig:results_ablation} that the full model achieves the lowest average MAE of 0.0438, the lowest MSE of 0.0032, and the highest $R^2$ score of 0.7734. These results reflect the model’s strong predictive accuracy and explanatory power.

Removing any single element leads to a decline in performance compared to the full model. Removing causality or correlation filtering increases error metrics, with the removal of causality showing the greatest effect, up to 0.2056 on average $R^2$. Disabling the causality mask also produces worse results, confirming the value of incorporating the causal mask into the attention layer. When two components are removed simultaneously, the model’s performance declines further. Configurations lacking the causality mask and another component yield the highest MAE and MSE values, along with an $R^2$ near 0 (see the purple and brown boxplots in Fig. \ref{fig:results_ablation}), performing even worse than the base model.

Overall, removing causality, correlation, or the transformer mechanism individually leads to a noticeable decline in the model's predictive accuracy and explanatory power. Furthermore, the simultaneous removal of multiple components results in a compounded decrease in performance, underscoring the synergistic contributions of causality, correlation, and the transformer mask to the model's effectiveness. %The base model demonstrated higher error metrics and a lower $R^2$ score, confirming that the inclusion of causality, correlation, and the transformer mask simultaneously enhances the model's performance compared to the base model.


\begin{table}[H]
\centering
\footnotesize
\resizebox{\textwidth}{!}{%
\begin{tabular}{l|c|c|c|c}
\toprule
 & \multicolumn{4}{c}{\textbf{Metric ($\pm$ Std Dev)}} \\ 
\cmidrule(lr){2-5}
\textbf{Model} & \textbf{MAE} & \textbf{MSE} & \textbf{R\textsuperscript{2}} & \textbf{Time (s)} \\ 
\midrule

Full Model
& \textbf{0.0438 $\pm$ 0.0143}
& \textbf{0.0032 $\pm$ 0.0019} 
& \textbf{0.7734 $\pm$ 0.1323} 
& \textbf{27.37 $\pm$ 0.7888} \\ 

Without Causality
& 0.0811 $\pm$ 0.0291 
& 0.0107 $\pm$ 0.00697 
& 0.2056 $\pm$ 0.4920 
& 28.07 $\pm$ 0.2746 \\ 

Without Correlation 
& 0.0794 $\pm$ 0.0095 
& 0.0097 $\pm$ 0.0028 
& 0.3088 $\pm$ 0.4920 
& 29.33 $\pm$ 0.2746 \\ 

Without Attention Mask
& 0.0601 $\pm$ 0.0194 
& 0.0060 $\pm$ 0.0041 
& 0.5600 $\pm$ 0.3056 
& 28.88 $\pm$ 0.9250 \\ 

Without Causality \& Correlation
& 0.0777 $\pm$ 0.0163 
& 0.0104 $\pm$ 0.0048 
& 0.1964 $\pm$ 0.4411 
& 29.21 $\pm$ 0.8247 \\ 

Without Causality \& Attention Mask
& 0.1067 $\pm$ 0.0360 
& 0.0167 $\pm$ 0.0100 
& -0.2059 $\pm$ 0.6280 
& 27.88 $\pm$ 0.9632 \\ 

Without Correlation \& Attention Mask
& 0.0933 $\pm$ 0.0283 
& 0.0149 $\pm$ 0.0076 
& -0.0916 $\pm$ 0.6744 
& 29.12 $\pm$ 0.6980 \\ 

Base\_model
& 0.0769 $\pm$ 0.0117 
& 0.0100 $\pm$ 0.0040 
& 0.2514 $\pm$ 0.2653 
& 29.21 $\pm$ 1.0849 \\ 

\bottomrule
\end{tabular}%
}
\caption{Cross-Validation Performance - Ablation Test}
\label{tab:results_ablation}
\end{table}



\begin{figure}[ht]
\centering
    \includegraphics[width = 0.85\linewidth]{Figures/ablation_r2.pdf}
    \caption{$R^2$ Score Distribution for different methods}
    \label{fig:results_ablation}
\end{figure}


\begin{figure}[H]
\centering
    \includegraphics[width = 0.9\linewidth]{Figures/ablation_train_loss.pdf}
    \caption{Train loss for ablation study}
    \label{fig:train_loss_ablation}
\end{figure}


\begin{figure}[H]
\centering
    \includegraphics[width = 0.9\linewidth]{Figures/ablation_val_loss.pdf}
    \caption{Validation loss for ablation study}
    \label{fig:validation_loss_ablation}
\end{figure}


 
\newpage 
{\color{black}
\subsection{Case Study: Generalizability of the Methodology} \label{subsec:urban_repro}

To validate the robustness and generalizability of our methodology, we applied the entire pipeline to a second case study: Salt Lake City, Utah (USA). For this study, we sourced traffic data from the Utah Department of Transportation (UDOT), which operates non-intrusive, side-fire radar stations providing vehicle counts \footnote{\url{https://www.udot.utah.gov/connect/business/traffic-data/}}. We selected four stations located on the main roads to the city (north, south, east, and west) and aggregated their data into weekly total counts to serve as the target variable $Y$. Figure \ref{fig:slc_region_map} shows the study area and the selected station locations.


\begin{figure}[ht]
\centering
    \includegraphics[width = 0.5\linewidth]{Figures/map_slc.pdf}
    \caption{Case Study of Salt Lake City}
    \label{fig:slc_region_map}
\end{figure}

To study visitation demand, we started with an initial set of five Google Trends terms specific to the new case study: ``Salt Lake City'', ``Weather in Salt Lake City'', ``Salt Lake City airport'', ``Great Salt Lake'', and ``Twin Peaks''. We then expanded this initial set using the PyTrends API, obtaining a total of 172 time series.

Next, we proceeded with the semantic validation of the expanded set. Since this corpus is in English, we used BERT (base, uncased) to ensure the tokenizer and embedding space were aligned with the data's linguistic context. We used the reference string ``Salt Lake City visitor attractions and travel'' for this validation step (see Equation \ref{eq:cosine_sim}). Following our methodology, we generated lagged series and filtered them using the Granger causality test (Equation \ref{eq:equation_causality}) and the correlation test (Equation \ref{eq:equation_correlation}).


We then trained the same set of baseline models and our GCT {\color{blue}architecture} under the identical 10-fold TimeSeriesSplit cross-validation scheme. The resulting metrics are reported in Table \ref{tab:slc_results} and Figure \ref{fig:urban_study}. As shown in Table \ref{tab:slc_results}, GCT again achieves the best performance across all metrics, obtaining the highest $R^2$ score ($0.5058$) and the lowest MAE ($0.0501$) and MSE ($0.0037$). This represents an $R^2$ improvement of 47.4\% over the next-best performing baseline, the Univariate LSTM with Attention ($R^2=0.3432$), and a 52.5\% improvement over TimesFM ($R^2=0.3317$). These results confirm the model's robustness and its ability to generalize effectively to the new context.
}


\begin{table}[ht]
\centering
\footnotesize
\resizebox{\textwidth}{!}{%
\begin{tabular}{l|c|c|c|c}
\toprule
 & \multicolumn{4}{c}{\textbf{Metric ($\pm$ Std Dev)}} \\ 
\cmidrule(lr){2-5}
\textbf{Model} & \textbf{MAE} & \textbf{MSE} & \textbf{R\textsuperscript{2}} & \textbf{Time (s)} \\ 
\midrule

ARIMA 
& 0.1540 $\pm$ 0.0215 
& 0.0250 $\pm$ 0.0723 
& -0.2425 $\pm$ 0.3322 
& 1.1639 $\pm$ 0.1419 \\ 

ARIMAX 
& 0.1343 $\pm$ 0.0246 
& 0.0186 $\pm$ 0.0054 
& 0.0147 $\pm$ 0.2260 
& \textbf{0.8832 $\pm$ 0.0374} \\ 

Univariate GRU
& 0.1120 $\pm$ 0.0623 
& 0.0168 $\pm$ 0.0110 
& 0.1563 $\pm$ 0.1755 
& 21.2521 $\pm$ 3.0565 \\

Multivariate GRU
& 0.1632 $\pm$ 0.0083 
& 0.0270 $\pm$ 0.0052 
& -0.2799 $\pm$ 0.3287 
& 24.9496 $\pm$ 2.6463 \\

Univariate LSTM
& 0.1530 $\pm$ 0.0199 
& 0.0245 $\pm$ 0.0044 
& -0.2415 $\pm$ 0.0569 
& 22.6211 $\pm$ 0.1947 \\

Univariate LSTM with Attention
& 0.0910 $\pm$ 0.0251 
& 0.0125 $\pm$ 0.0062 
& 0.3432 $\pm$ 0.2110
& 30.5435 $\pm$ 1.1426 \\

Multivariate LSTM
& 0.1410 $\pm$ 0.0259 
& 0.0215 $\pm$ 0.0083 
& -0.0802 $\pm$ 0.3735 
& 30.4370 $\pm$ 0.6659 \\

Multivariate Conv1d+LSTM \cite{long2021air}
& 0.1680 $\pm$ 0.0195 
& 0.0285 $\pm$ 0.0041 
& -0.3340 $\pm$ 0.3562 
& 26.9420 $\pm$ 0.3810 \\

Moirai \cite{woo2024unified}
& 0.1350 $\pm$ 0.0262
& 0.0190 $\pm$ 0.0080
& 0.0096 $\pm$ 0.2385
& 5.6800 $\pm$ 0.0929 \\

TimesFM \cite{das2024decoder}
& 0.0932 $\pm$ 0.0118
& 0.0133 $\pm$ 0.0051
& 0.3317 $\pm$ 0.1712
& 2.5020 $\pm$ 0.4969 \\

cLSTM \cite{tank2021neural}
& 0.1450 $\pm$ 0.0401 
& 0.0225 $\pm$ 0.0235 
& -0.1215 $\pm$ 0.2365 
& 13.3290 $\pm$ 0.6383 \\

cMLP \cite{tank2021neural}
& 0.1480 $\pm$ 0.0180 
& 0.0230 $\pm$ 0.0057 
& -0.1481 $\pm$ 0.1660 
& 2.2050 $\pm$ 0.6029 \\

Causalformer \cite{kong2024causalformer}
& 0.1050 $\pm$ 0.0127 
& 0.0155 $\pm$ 0.0027 
& 0.2261 $\pm$ 0.2108 
& 9.3530 $\pm$ 0.7064 \\

GCT (Our model)
& \textbf{0.0501 $\pm$ 0.0164}
& \textbf{0.0037 $\pm$ 0.0022}
& \textbf{0.5058 $\pm$ 0.1327}
& 36.9810 $\pm$ 2.2605 \\
\bottomrule
\end{tabular}%
}
\caption{Cross-Validation Performance - Methods with Execution Times - Salt Lake City}
\label{tab:slc_results}
\end{table}


\begin{figure}[ht]
\centering
    \includegraphics[width = 1\linewidth]{Figures/r2_boxplot_vertical_utah.pdf}
    \caption{Comparison of $R^2$ results - Salt Lake City}
    \label{fig:urban_study}
\end{figure}



\subsection{Discussion}\label{sec:discussion}


The comparison across statistical and deep learning models, both in their univariate and multivariate cases, provides insights into their capabilities and limitations. As summarized in {\color{black} our first case study} (Table \ref{tab:sota}), statistical models (ARIMA, ARIMAX) performed poorly, with low or even negative $R^2$ scores, likely due to their linear nature and inability to capture complex, non-linear patterns in the data. Univariate deep learning models, such as GRU, LSTM, and LSTM with Attention, showed improvements with $R^2$ scores between $0.3847$ and $0.4596$. However, their limited capacity to incorporate contextual information constrains their explanatory power. Multivariate GRU and LSTM models achieved marginally better scores but remained low, see Fig. \ref{fig:sota}, indicating their limitations in capturing the dataset's full complexity. Causality-aware methods, such as cLSTM and cMLP, also fail to adequately fit these data, as evidenced by their poor average $R^2$ values ($-0.7696$ for cLSTM and $0.0455$ for cMLP). These architectures lack attention mechanisms and consequently cannot capture the complex temporal relationships present in real time series, limiting their predictive capacity. Causalformer, which incorporates attention, proves insufficient for these data, achieving an average $R^2$ of $0.289$. CausalFormer accepts input from a single predefined observation window and cannot identify multi-scale causal patterns \cite{kong2024causalformer}, which is a significant limitation in this analysis. Our proposed model surpasses all models across all metrics. The lower MAE ($0.0438$), lower MSE ($0.0032$), and higher $R^2$ ($0.7734$) than other solutions indicate a better fit to the data, due to the model's ability to integrate and effectively capture both contextual and causal relationships. This adaptability allows it to explain a higher variance in the data, providing more reliable predictions. Despite its longer execution time (averaging $27.37$ seconds compared to other algorithms, which required at most $23.63$ seconds), the improvement in performance justifies this trade-off.

{\color{black} To test the method's robustness, we further evaluated its generalizability in the second context of Salt Lake City (Table \ref{tab:slc_results}). This scenario proved more challenging, as most baseline models, including standard multivariate deep learning approaches, yielded negative $R^2$ scores. The best-performing baselines were the large foundation models, TimesFM ($R^2=0.3317$), and the Univariate LSTM with Attention ($R^2=0.3432$). Once again, our proposed GCT {\color{blue}architecture} significantly outperformed all other methods, achieving the highest $R^2$ score ($0.5058$) and the lowest errors (MAE $0.0501$, MSE $0.0037$). This represents a 47.4\% increase in $R^2$ compared to the next-best baseline. The consistent success of GCT in both a sparse, rural setting (Alpujarra, Spain) and a dense setting (Salt Lake City, USA) highlights the robustness of our methodology.}






{\color{black}
\subsection{Limitations and Future Work} \label{sec:limitations}

Although the results provide a superior performance with respect to the baseline methods, the proposed methodology has some limitations. First, the multi-step preprocessing and the complexity of GCT increase training time, which can cause delays when using large datasets or real-time applications. Second, the use of a Granger-based linear screening method serves only as a filtering step to select relevant inputs. The retained signals are then modeled by the non-linear GCT architecture (LSTM and Transformer), which is designed to capture complex, non-linear interactions. Performance in scenarios dominated by purely non-linear relationships may still require careful lag selection and more extensive hyperparameter tuning. Third, even though leveraging Google Trends offers a valuable, publicly available proxy for behavioral dynamics, GCT reliance on this source limits portability to regions with low Google adoption (e.g., China, Russia).

Additionally, with respect to the data used, our methodology could be extended to support other datasets, such as forecasting, land use features, road network features, which can be helpful to improve the prediction accuracy. Unfortunately, for our rural case study, we could not find reliable datasets. The region lacks a local meteorological station, and triangulated estimates in mountains with microclimate proved too noisy for reliable integration. Similarly, there are no open datasets with sufficient spatial/temporal resolution for land use or network topology in the area. Given these constraints, we prioritized publicly accessible Google Trends time series and did not include these other external datasets in our methodology. Another limitation stems from the destination-wide scope of our methodology. The model is designed to forecast the total aggregated weekly vehicle counts to a study area, aiming to capture destination-wide level demand, because normally tourist go to different spots in the area. Consequently, it does not disaggregate or model the fine-grained mobility within that area, such as per-location forecasts. {\color{blue}Therefore, a direct comparison of model performance across individual sub-locations or villages within the aggregated study area is not possible with the current destination-wide approach.}


As future work, we aim to optimize preprocessing and model design to improve scalability and real-time performance. Working with the Alpujara region council, we plan to install additional sensors (e.g., a local weather station) to support multimodal sensor fusion. Finally, we will extend the methodology to support fine-grained, per-location forecasts, which would be especially useful in areas with distinct intra-destination spots.
}










\section{Conclusion} \label{sec:conclusions}

This study presents {\color{blue}GCT}, a transformer-based architecture that incorporates a causal mask to modify the attention similarity weight matrix. This causal mask uses Granger Causality tests to identify causal relationships among inputs in a multivariate NN, adjusting the attention weights accordingly. By doing so, the GCT ensures that the model focuses not only on the temporal similarities captured by the attention mechanism but also on the causal relationships among inputs.


To validate this proposal, we designed a methodology that includes an initial feature validation step using a language model, followed by a feature selection step based on causal and correlation filters, and finally a prediction model using our GCT architecture. {\color{black} We first applied this methodology to a case study for predicting weekly vehicle counts in a rural area.} The experimental results demonstrate that the proposed approach outperforms other statistical and deep learning models, achieving improvements of up to 30\% in MAE, 56\% in MSE, and 68\% in $R^2$ compared to the best performing baseline method. {\color{black} Furthermore, we confirmed the method's generalizability in a second case study (Salt Lake City), where GCT again significantly outperformed all baselines, achieving a 47\% increase in $R^2$ over the strongest competitor. The consistent success in both contexts validates the robustness and wide applicability of our approach.}


%%
%% The acknowledgments section is defined using the "acks" environment
%% (and NOT an unnumbered section). This ensures the proper
%% identification of the section in the article metadata, and the
%% consistent spelling of the heading.
\begin{acks}
This work was supported in part by Grant C-SEJ-128-UGR23 funded by Consejería de Universidad, Investigación e Innovación and by ERDF Andalusia Program 2021-2027; project PID2023-149185OBI00 funded by MICIU/AEI/10.13039/501100011033 and co-funded by ERDF/EU.
\end{acks}

%%
%% The next two lines define the bibliography style to be used, and
%% the bibliography file.
\bibliographystyle{ACM-Reference-Format}
\bibliography{sample-base}


%%
%% If your work has an appendix, this is the place to put it.
%\appendix


\end{document}
\endinput
%%
%% End of file `sample-acmlarge.tex'.
